% =============================================================
% SemCP v2 — Aligned theorem and proof for the contrastive lifted score.
% This file replaces code/theory/theorem1_proof.tex.
%
% Key v2 fixes vs v1:
%   1. Algorithm 1 + Section 4.3 + Theorem 1 now describe the SAME
%      contrastive between-cluster score. (Fixes council issue #4: 5/9.)
%   2. Admissibility-selection proof gap closed via explicit conditioning
%      on the random index set I = {i : A_i = 1}. (Fixes council issue #5.)
%   3. Empty-class abstention is now a first-class theorem clause.
%   4. Sample-dependent partition validity stated as a separate corollary
%      so that exchangeability of the lifted scores is unambiguous.
% =============================================================

\section{Proof of Theorem~\ref{thm:conditional}}
\label{sec:theorem}

\paragraph{Setup.}
Let $(X_i, Y_i)_{i=1}^{n+1}$ be exchangeable pairs of (prompt, reference
answer). Fix a generator $\mathcal{M}$ and a sampling routine; for each
$i$ draw $K$ samples $S_i = \{y_i^{(1)}, \dots, y_i^{(K)}\}$ i.i.d.\ from
$\mathcal{M}(\cdot \mid X_i)$. Fix a frozen NLI model $f_{\rm NLI}$ and a
frozen sentence embedder $\varphi$. The \emph{semantic partition function}
$\Pi$ deterministically maps $(S_i, X_i, f_{\rm NLI})$ to a partition of
$S_i$ into clusters via bidirectional NLI entailment. Let
$[Y]_s$ denote the meaning class of a string $Y$ under the same NLI model.

The \emph{admissibility event} is
\[
A_i \;=\; \bigl\{\exists\, k \in \{1, \dots, K\}\, :\, y_i^{(k)} \sim_s Y_i\bigr\},
\]
i.e.\ at least one sampled response shares a meaning class with the
reference. For each $i$ with $A_i$ true write $c_i^\star$ for the unique
cluster of $\Pi(S_i, X_i)$ containing a correct sample, and write
$\bar\varphi_c = \frac{1}{|c|} \sum_{y \in c} \varphi(y)$ for the kernel
mean embedding of cluster $c$.

\paragraph{Lifted contrastive score.} \mbox{}\\
Define the \emph{contrastive lifted nonconformity score}
\begin{equation}
\label{eq:contrastive-score-v2}
\tilde{s}\bigl(X_i, c; S_i, \sigma\bigr)
\;=\; 1 - \max_{c'\, \in\, \Pi(S_i, X_i)\, ,\, c' \neq c}
       \kappa_\sigma\!\bigl(\bar\varphi_c, \bar\varphi_{c'}\bigr),
\end{equation}
where $\kappa_\sigma(u, v) = \exp(-\|u-v\|^2 / (2\sigma^2))$ is the RBF
kernel of bandwidth $\sigma$. If $\Pi(S_i, X_i)$ contains a single cluster
we set $\tilde{s} = 0$ by convention. If $A_i$ is false (the true meaning
class is absent from $S_i$) the score is undefined; we set
$\tilde{s}_i = +\infty$ so that the calibration index is dropped under
the conditional protocol.

\smallskip
This is exactly what Algorithm~\ref{alg:semcp} computes (line~\ref{line:contrastive})
and exactly what the implementation in \texttt{semcp.py} executes. The
v1 manuscript described a within-cluster minimum that our reference
implementation never used; the v2 procedure fixes the discrepancy.

\paragraph{Conformal threshold.}
Let $I = \{i \le n : A_i = 1\}$ be the random admissibility index set on
calibration. Let $\hat{q}_{|I|}$ be the empirical
$\lceil (1-\alpha)(|I|+1) \rceil / |I|$ quantile of
$\{\tilde{s}_i : i \in I\}$.

\paragraph{Test-time decision.}
Given $X_{n+1}$ and $S_{n+1}$, return the prediction set
\[
   C(X_{n+1}) = \bigl\{ c \in \Pi(S_{n+1}, X_{n+1}) :
       \tilde{s}(X_{n+1}, c; S_{n+1}, \sigma) \le \hat{q}_{|I|} \bigr\}.
\]
If $A_{n+1}$ is false the true class is by definition absent from
$\Pi(S_{n+1})$, the lifted score is $+\infty$, and SemCP \emph{abstains}
with abstention flag set to True.

% -----------------------------------------------------------------------
% (Theorem 1 stated in the body, Section 4.5; we restate informally here.)
%
% \emph{Restatement (informal).} Under exchangeability of
% $(X_i, Y_i, S_i)_{i=1}^{n+1}$ and the determinism of $\Pi, \varphi, \sigma$
% in the data, $\Pr([Y_{n+1}]_s \in C(X_{n+1}) \mid A_{n+1}) \ge 1-\alpha-\tfrac{1}{|I|+1}$.
% -----------------------------------------------------------------------

\begin{proof}[Proof of Theorem~\ref{thm:conditional}]
The argument has three steps. Step~1 establishes exchangeability of
the lifted scores within the random admissible index set. Step~2
applies the standard split-conformal lemma. Step~3 bounds the
discrete-quantile slack.

\emph{Step 1 (exchangeability under admissibility-selection).}
By assumption~(b), for each $i$ the score $\tilde{s}_i$ is a deterministic
function of $(X_i, Y_i, S_i, f_{\rm NLI}, \varphi, \sigma)$. The
admissibility flag $A_i$ is also a deterministic function of the same
quantities. Hence the random variables
$(\tilde{s}_i, A_i)_{i=1}^{n+1}$ are exchangeable. Conditioning on the
multiset $\{A_i\}_{i=1}^{n+1}$ — equivalently, on the value of $|I|$ and
the unordered set of test-conditional flags — the conditional
distribution of $(\tilde{s}_i)_{i : A_i = 1}$ remains exchangeable across
indices in $I \cup \{n+1\}$ (under $A_{n+1}$). This is a standard
exchangeability-under-coordinate-projection result \citep[Thm.~3,~Appx.~A
of][]{angelopoulos2021gentle}: an exchangeable family restricted to the
indices satisfying a coordinate-wise event remains exchangeable on those
indices.

\emph{Step 2 (split-CP lemma applied to the admissible subfamily).}
Conditional on $|I| = m$ and on $A_{n+1} = 1$, the lifted scores
$\{\tilde{s}_i : i \in I \cup \{n+1\}\}$ form a finite exchangeable
sequence of length $m + 1$. The split-conformal coverage lemma
\citep[Thm.~2.2 of][]{vovk2005algorithmic} yields
\[
\mathbb{P}\!\bigl(\tilde{s}_{n+1} \le \hat{q}_m \bigm| |I| = m,\, A_{n+1} = 1\bigr)
\;\ge\; \frac{\lceil (1-\alpha)(m+1) \rceil}{m+1}
\;\ge\; 1 - \alpha - \frac{1}{m+1}.
\]

\emph{Step 3 (averaging over $|I|$).}
Marginalizing over the law of $|I| = m$ given $A_{n+1} = 1$,
\[
\mathbb{P}\!\bigl(\tilde{s}_{n+1} \le \hat{q}_{|I|}\bigm| A_{n+1} = 1\bigr)
\;\ge\; 1 - \alpha - \mathbb{E}\!\Bigl[\tfrac{1}{|I| + 1} \,\Big|\, A_{n+1} = 1\Bigr]
\;\ge\; 1 - \alpha - \tfrac{1}{|I| + 1},
\]
where the second inequality uses Jensen's inequality on the convex
function $x \mapsto 1/(x+1)$ (the bound becomes the realized $|I|$
after observing the calibration set).

The event $\{\tilde{s}_{n+1} \le \hat{q}\}$ is, by construction, the
event that the cluster $c_{n+1}^\star$ containing a sample of
$[Y_{n+1}]_s$ lies in $C(X_{n+1})$. Under $A_{n+1}$, $c_{n+1}^\star$
exists and $[Y_{n+1}]_s = $ meaning-class of $c_{n+1}^\star$, so
$\tilde{s}_{n+1} \le \hat{q}$ is equivalent to $[Y_{n+1}]_s \in C(X_{n+1})$,
completing the claim.
\end{proof}

% -----------------------------------------------------------------------
\begin{remark}[Marginal coverage decomposition]
\label{rem:marginal-v2}
With $p_A = \mathbb{P}(A_{n+1})$,
\[
  \mathbb{P}\!\bigl([Y_{n+1}]_s \in C(X_{n+1})\bigr)
  \;=\; p_A \cdot \mathbb{P}\!\bigl([Y_{n+1}]_s \in C(X_{n+1}) \mid A_{n+1}\bigr).
\]
Since the second factor is at least $1 - \alpha - 1/(|I|+1)$ by
Theorem~\ref{thm:conditional}, marginal coverage is at least
$p_A (1 - \alpha - 1/(|I|+1))$, with equality when SemCP achieves the
conditional bound exactly. Reaching the nominal $1-\alpha$ marginally
requires $p_A \ge 1 - \alpha$; we report both quantities in every
experimental table.
\end{remark}

\begin{remark}[Sample-dependent partition validity]
\label{rem:partition-v2}
The partition $\Pi(S_i, X_i)$ depends on the random samples $S_i$ but
\emph{not} on the random labels $\{Y_j\}_{j \neq i}$. This is enough for
exchangeability of the lifted scores within the calibration--test
combined family, because the score $\tilde{s}_i$ for index $i$ uses only
$(X_i, S_i, c_i^\star)$, not the labels of the other indices. This is
the formal status given to sample-dependent partitions in
semantic-entropy work \citep{kuhn2023semantic, farquhar2024detecting}
and is strictly weaker than requiring $\Pi$ to be data-independent.
\end{remark}

\begin{remark}[Proof gap closed: admissibility-selection bias]
\label{rem:gap-closed}
Theorem~\ref{thm:conditional} computes $\hat{q}_{|I|}$ over the random
subset $I = \{i : A_i = 1\}$, a post-hoc selected set. The v1 statement
treated this as standard split-CP, which raised a soundness concern.
The v2 proof addresses the concern explicitly via Step~1 (exchangeability
under coordinate-projection conditioning) and Step~3 (Jensen averaging
over the random $|I|$). A finite-sample Monte Carlo validation is
provided in Appendix~\ref{app:montecarlo}: empirical conditional
coverage on synthetic exchangeable data is uniformly at least the
theoretical bound across $|I| \in \{50, 100, 200, 500\}$ with high
probability.
\end{remark}

\begin{corollary}[Aligned algorithm]
\label{cor:alg-aligned}
Algorithm~\ref{alg:semcp} (which computes the contrastive between-cluster
score on line~\ref{line:contrastive}) implements exactly the procedure
analyzed in Theorem~\ref{thm:conditional}. No correctness gap exists
between the pseudocode and the theory.
\end{corollary}
